\documentclass[a4paper]{article}
\usepackage{amsfonts}
\usepackage{amsmath}
\usepackage{amssymb}
\usepackage{graphicx}     

\begin{document}
  
\noindent \large Auteur: Stan Van Campenhout \\
\noindent \large Studierichting:  Informatica\\
\noindent \large Oefeningenreeks 6A nr. 10.3\\

\medskip

\normalsize

$ a, b, c, d $\\

Meetkundige rij met $ a + b = 28, c + d = 252 $\\

$ \Rightarrow a + aq = 28, aq^2 + aq^3 = 252 \Leftrightarrow a\left(1+q\right) = 28, a\left(q^2+q^3\right) = 252 $\\

$ \Rightarrow \dfrac{28}{1+q} = \dfrac{252}{q^2+q^3} \Leftrightarrow \left(1+q\right)252 = \left(q^2+q^3\right)252 $\\

$ \Leftrightarrow 28q^3+28q^2-252q-252 = 0 $\\

 \[ \begin{array}{c|rrrr}     
& 28 & 28 & -252 & -252 \\      
-1 & \downarrow  & -28 & 0 & 252   \\ \hline    
 & 28  & 0  &  \multicolumn{1}{c|}{-252}  & 0   
\end{array} \] \\

$ \Rightarrow \left(q+1\right)\left(28q^2-252\right) = 0 \Rightarrow  \left(q+1\right)\left(q-3\right)\left(q+3\right) = 0 $\\

Schrap $q = -1$ aangezien we deelden door $1+q $\\

$ \Rightarrow a \pm 3 = 28 \Leftrightarrow a \in \{7, -14\} $\\

$ \Rightarrow -14,42,-126,378$ en $ 7, 21, 63, 189 $\\

\begin{thebibliography}{999}
%\bibitem{a} Referentie
\end{thebibliography}
\end{document} 
